\section{Pacuit's neighborhood constructions: full proof interface}

\begin{annexitem}{annex:p8.pn.extensional}{thm:p08-pn-extensional}{Extensional truth, the dual translation and monotonic reading}
\begin{proof}
Write \(T(\alpha)=[![\alpha]!]_{\mathcal M}\). If
\(T(\alpha)=T(\beta)\), the two sets have the same membership
in every \(N(w)\), so \(T(\Box\alpha)=T(\Box\beta)\).
Validity of a biconditional throughout the model gives precisely
the required equality; its truth at one world does not. For example,
on \(W=\{a,b\}\), set \(N(a)=N(b)=\{\{a\}\}\),
\(V(p)=\{a\}\), and \(V(q)=W\). At a, \(p\leftrightarrow q\)
is true but \(\Box p\leftrightarrow\Box q\) is false.

By the defined diamond and classical complement,
\(\mathcal M,w\models\neg\Diamond\neg\alpha\) iff
\(W\setminus T(\neg\alpha)=T(\alpha)\) belongs to \(N(w)\).
This is the box clause, including when the truth set is empty.
For \(\langle\ ]\), a witnessing neighborhood contained in
\(T(\alpha)\) is also contained in any larger \(T(\beta)\).
On a supplemented family it is a witness exactly when
\(T(\alpha)\in N(w)\): supplementation gives one direction,
and the truth set itself is a witness for the other. In the example
above \(T(p)\subseteq T(q)\) but \(\Box p\) is true and
\(\Box q\) false, so exact-membership box is not generally monotonic.

For the source language with primitive diamond, let t fix atoms,
constants and box, commute with Boolean connectives, and send
\(\Diamond\alpha\) to \(\neg\Box\neg t(\alpha)\).
In E, Dual at \(\neg\alpha\), double-negation replacement under
box, and Boolean negation give
\(\Diamond\alpha\leftrightarrow\neg\Box\neg\alpha\).
Induction on a formula now proves \(\alpha\leftrightarrow t(\alpha)\):
Boolean cases use tautologies, the box case RE, and the diamond case
this derived identity and the inner induction equivalence. The same
induction gives semantic truth preservation. Translated Dual is a
theorem of the defined-diamond language; translated MP and RE are
valid theorem inferences. Conversely each defined-diamond proof is
a source-language proof after writing its expansions. Thus this is
a conservative theorem translation, not a change of modal semantics.
\end{proof}
\end{annexitem}

\begin{annexitem}{annex:p8.pn.correspondence}{thm:p08-pn-correspondence}{The separate all-valuations correspondences}
\begin{proof}
For M, \(T(\alpha\land\beta)=T(\alpha)\cap T(\beta)\)
is contained in both factors. Supplementation therefore makes
its membership imply both memberships. Conversely, a failure of
supplementation supplies \(X\subseteq Y\), \(X\in N(w)\),
\(Y\notin N(w)\). Set \(V(p)=X\), \(V(q)=Y\).
Then \(T(p\land q)=X\), so the M instance at w is false.
For C, intersection closure turns the two memberships into that
of the intersection. A failure supplies \(X,Y\in N(w)\)
with \(X\cap Y\notin N(w)\); the valuation
\(V(p)=X,V(q)=Y\) falsifies its C instance. These arbitrary-set
converses use the two declared atoms. Since \(T(\topv)=W\),
N is valid exactly when the unit belongs to every family.

For D, \(\Box\alpha\to\Diamond\alpha\), the diamond
clause says that membership of \(T(\alpha)\) must exclude
membership of its complement. Conversely any complementary pair
is witnessed by \(V(p)=X\). Finally
\(\Diamond\topv\) is true exactly when \(\varnothing\notin N(w)\).
In a regular family a complementary pair generates the empty set
by intersection; and if the empty set belongs, supplementation
also inserts W, a complementary pair. Thus the two conditions
are equivalent for regular families, including the empty family.
They need not agree for merely monotonic ones. On
\(W=\{a,b\}\), take at every world
\(N(w)=\{\{a\},\{b\},W\}\). This is supplemented and
excludes the empty set, but contains a complementary pair.
With \(V(p)=\{a\}\), \(\Diamond\topv\) is true and
\(\Box p\to\Diamond p\) false at every world.
\end{proof}
\end{annexitem}

\begin{annexitem}{annex:p8.pn.derived-rules}{thm:p08-pn-derived-rules}{All source rule identities, including the zero-antecedent boundary}
\begin{proof}
First suppose \(\vdash_L\alpha\leftrightarrow\beta\).
At an unchanged atomic or constant occurrence use reflexivity;
at a replaced occurrence use the assumed equivalence. Induct up
the formula tree. Negation preserves equivalence by a tautology,
and conjunction preserves the two induction equivalences by a
tautology and MP. At box use RE on the inner theorem equivalence.
Diamond follows by its expansion. Multiple chosen occurrences
are handled by this same induction, or by composing finitely many
single-occurrence replacements. No equivalence depends on a local premise.

If \(\vdash_{EM}\alpha\to\beta\), propositional reasoning
gives \(\vdash_{EM}\alpha\leftrightarrow(\alpha\land\beta)\).
RE and M imply \(\vdash_{EM}\Box\alpha\to\Box\beta\).
Conversely RM applied to the two tautological projections of
\(\alpha\land\beta\) proves M. Hence the least theorem sets
E+M and E+RM agree. Nec applied to \(\topv\) proves N.
If \(\vdash_{EN}\alpha\), then
\(\vdash_{EN}\topv\leftrightarrow\alpha\); RE and N prove
\(\Box\alpha\). This establishes EN=E+Nec.

In EMC, RM boxes the tautology
\(((\alpha\to\beta)\land\alpha)\to\beta\).
C combines \(\Box(\alpha\to\beta)\) with \(\Box\alpha\)
into its boxed antecedent. MP and a propositional currying
tautology give K. In normal K, Nec and K first give RM: necessitate
the theorem \(\alpha\to\beta\) and use K. The two projections
then give M. To obtain C, necessitate
\(\alpha\to(\beta\to(\alpha\land\beta))\); two K instances
give \(\Box\alpha\to(\Box\beta\to\Box(\alpha\land\beta))\).
Nec also gives N. Finally RM applied to both directions of a
theorem biconditional yields RE, so the E component is included.
The two inclusions prove EMCN=K. C and M commute as axiom additions,
so the source name ECM denotes EMC.

For any classical L containing M and C, assume the theorem
\((\alpha_1\land\cdots\land\alpha_k)\to\beta\), with
\(k\geq1\). Repeated C derives
\((\Box\alpha_1\land\cdots\land\Box\alpha_k)
\to\Box(\alpha_1\land\cdots\land\alpha_k)\);
RM gives its implication to \(\Box\beta\). For k=1 no C
step is needed. Conversely the k=1 case gives RM and hence M;
the k=2 case applied to the identity conjunction gives C.
The k=0 case means precisely that theorem \(\beta\) yields
theorem \(\Box\beta\), so adding it gives Nec and N. This
proves the regular and normal characterizations with their exact arities.

To connect this theorem calculus to premise derivations, fix a
finite list of used premises with conjunction \(\delta\).
Induct over a derivation in which modal inferences have independent
theorem proofs. A premise line has theorem \(\delta\to\alpha\);
a theorem line has that implication by propositional weakening;
and MP combines the two induction implications using its Boolean
tautology. Every modal line is a theorem line. Thus a premise
derivation yields \(\vdash_L\delta\to\varphi\).
Conversely derive the finite conjunction from its premises by
propositional tautologies and MP, then use the theorem implication.
The empty list gives theoremhood in both directions. The established
theorem-set identities therefore preserve precisely this local deduction.
\end{proof}
\end{annexitem}

\begin{annexitem}{annex:p8.pn.derived-principles}{thm:p08-pn-derived-principles}{Each monotonic and regular source consequence}
\begin{proof}
The two E equivalences follow directly from the diamond definition
and replacement of \(\neg\topv\) by \(\bottom\), or
\(\neg\bottom\) by \(\topv\), under box. For the monotonic
principles, write \(A=T(\varphi)\), \(B=T(\psi)\).
Both B and \(W\setminus A\) are subsets of
\((W\setminus A)\cup B=T(\varphi\to\psi)\), so their
membership implies its membership. If
\(\Diamond(\varphi\to\psi)\) is false, then
\(A\setminus B=T(\neg(\varphi\to\psi))\) is accepted.
It is contained in \((W\setminus B)\cup A\), so
\(\Box(\psi\to\varphi)\) is true. If either
\(\Diamond\varphi\) or \(\Diamond\psi\) is false,
its accepted complement is contained in
\(W\setminus(A\cap B)\); thus \(\Diamond(\varphi\land\psi)\)
is false. This proves its displayed implication by contraposition.

Theorem diamond monotonicity is also derivable directly: from
theorem \(\varphi\to\psi\) derive theorem
\(\neg\psi\to\neg\varphi\), apply boxed RM, and
contrapose the result, expanding the defined diamonds. Every
inference has a theorem premise.

For the first regular principle suppose
\(I=(W\setminus A)\cup B\) is accepted but
\(\Diamond\psi\) false, so \(W\setminus B\) is accepted.
Intersection closure accepts \(I\cap(W\setminus B)\),
a subset of \(W\setminus A\). Supplementation accepts
\(W\setminus A\), making \(\Diamond\varphi\) false.
For the boxed biconditional its truth set is contained in both
implication truth sets, so M accepts them. If A is accepted,
C accepts its intersection with \((W\setminus A)\cup B\),
which is \(A\cap B\); M then accepts B. Interchange A and B
for the converse boxed implication. Finally if A is accepted
but \(\Diamond(\varphi\land\psi)\) false, intersection
closure accepts \(A\cap(W\setminus(A\cap B))=A\setminus B\).
Supplementation accepts \(W\setminus B\), making
\(\Diamond\psi\) false. This proves the third regular principle.

If \(\Diamond\topv\) is false, the empty set is accepted,
and supplementation accepts every set. Both \(\Box\varphi\)
and \(\Box\neg\varphi\) are then true, so
\(\Box\varphi\to\Diamond\varphi\) is false. This proves
the monotonic implication to \(\Diamond\topv\). In a regular
family, truth of \(\Diamond\topv\) excludes complementary
accepted pairs by the correspondence proof, so implies D.
The two-element supplemented family excluding the empty set
but containing its two complementary singletons witnesses failure
of this converse without C. All the displayed modal schemas have
now been proved valid on the appropriate EM or EMC frame class.
Their full completeness proofs give EM or EMC theoremhood,
respectively; every monotonic or regular extension contains those
base theorem sets. Thus the asserted general theorem status follows,
without presupposing normality or adding N.
\end{proof}
\end{annexitem}

\begin{annexitem}{annex:p8.pn.finite-separations}{thm:p08-pn-finite-separations}{Every world and valuation of the separating frames}
\begin{proof}
First retain the source's explicit two-world C counterexample:
\(W=\{a,b\}\), \(N(a)=\{\{a\},\{b\}\}\),
\(N(b)=\{\varnothing\}\), \(V(p)=\{a\}\), and
\(V(q)=\{b\}\). At a both boxes are true, while
\(T(p\land q)=\varnothing\notin N(a)\), so their
conjunct's box is false. The family at b is part of this
original example; it is not left unspecified or normalized.
The following stronger separating frames additionally retain
the two other closure conditions in each selected failure.

Each following family is assigned to \emph{every} world of its
specified nonempty domain, so no unspecified world affects frame
validity. To retain C and N while failing M use
\(W=\{a,b,c\}\) and \(\mathcal U=\{\{a\},W\}\).
It is intersection-closed and unit-containing, but
\(\{a\}\subseteq\{a,b\}\) violates supplementation. With
\(V(p)=\{a\},V(q)=\{a,b\}\),
\(\Box(p\land q)\) is true and \(\Box q\) false.
To retain M and N while failing C use
\(W=\{a,b\}\), \(\mathcal U=\{\{a\},\{b\},W\}\).
It is supplemented and unit-containing; with
\(V(p)=\{a\},V(q)=\{b\}\), both boxes are true and
their conjunct's box false. To retain M and C while failing N,
use \(W=\{a\}\), \(\mathcal U=\varnothing\).
The two closure implications are vacuously true while
\(\Box\topv\) is false. These prove independent schema
failures by soundness of the corresponding theorem rules.

The first, unit-containing, three-world example also fails K:
take \(V(p)=\{a\}\), \(V(q)=\{a,b\}\).
The truth set of \(p\to q\) is W; its box and \(\Box p\)
are true while \(\Box q\) is false. Thus EN is not K.
The empty-family example validates EMC but fails N, separating
regular from normal logic.

For the source's EK example put \(W=\{a,b,c,d\}\) and
\(N(x)=\{\{a,b\},\{a,c\}\}\) for every x. If
\(T(\alpha)=\{a,b\}\), then \(T(\alpha\to\beta)\)
contains c and d, so is neither listed neighborhood. If
\(T(\alpha)=\{a,c\}\), that implication contains b and d
and again is neither. Therefore the antecedents
\(\Box(\alpha\to\beta)\) and \(\Box\alpha\) cannot
both be true. Every K instance is valid under every valuation.
Take \(V(p)=\{a,b\},V(q)=\{a,c\}\).
Then \(\Box p\land\Box q\) is true, but
\(T(p\land q)=\{a\}\) is not a neighborhood, so C fails.
Also \(T((p\lor q)\land p)=\{a,b\}\) is a neighborhood
while \(T(p\lor q)=\{a,b,c\}\) is not, so M fails.
N fails because W is absent. Since PC, Dual, RE and MP are
sound on this frame, M and C are not theorems of EK.

For premise discipline take \(W=\{a,b\}\),
\(N(a)=N(b)=\{W\}\), and \(V(p)=\{a\}\).
At a, p is true and \(\Box p\) false although the frame has
the unit. On any unit model where p is true at all worlds its
truth set is W, hence \(\Box p\) is true at all worlds. For
frame-global consequence with premise p, no nonempty frame
validates p under all valuations (set \(V(p)=\varnothing\)),
so the implication is instead vacuous. These are distinct reasons
for two distinct global-consequence relations.
\end{proof}
\end{annexitem}

\begin{annexitem}{annex:p8.pn.proof-set-lemmas}{thm:p08-pn-proof-set-lemmas}{Lindenbaum extension and all Boolean and entailment claims}
\begin{proof}
Enumerate the countable language as \(\theta_0,\theta_1,\ldots\).
Begin with any L-consistent \(\Gamma_0\), even if infinite.
At stage n adjoin \(\theta_n\) if consistency is preserved;
otherwise adjoin \(\neg\theta_n\). The latter choice preserves
consistency. Otherwise finite conjunctions \(\delta,\epsilon\)
from \(\Gamma_n\) would give theorems
\((\delta\land\theta_n)\to\bottom\) and
\((\epsilon\land\neg\theta_n)\to\bottom\).
Classical reasoning would give
\((\delta\land\epsilon)\to\bottom\), contradicting
consistency of \(\Gamma_n\). Let \(\Delta=\bigcup_n\Gamma_n\).
A finite inconsistent subset would occur at one stage, so
\(\Delta\) is consistent. It decides each formula and is maximal:
adjoining any omitted \(\theta\) conflicts with its existing negation.

A maximal theory contains every theorem: omitting one would
insert its negation and give a finite contradiction. It cannot
contain both \(\alpha\) and \(\neg\alpha\); by decision,
exactly one belongs. If \(\alpha,\alpha\to\beta\) belong
but \(\beta\) does not, the three formulas with
\(\neg\beta\) form an inconsistent finite set. Thus it is
MP-closed. The tautological projections and introduction of
conjunction now give
\(\alpha\land\beta\in\Delta\) iff both factors belong.
The negation and conjunction facts give the disjunction truth
table, and the other classical connectives follow by expansion.
Every theory contains \(\topv\) and none contains \(\bottom\).
Taking world sets proves the asserted intersection, complement,
union, unit and empty proof-set identities.

A theorem \(\alpha\to\beta\) gives
\(|\alpha|_L\subseteq|\beta|_L\) by MP closure. If that
implication is not a theorem, \(\{\alpha,\neg\beta\}\)
is consistent: inconsistency of this pair would, by a propositional
tautology, prove the implication. Extend the pair by Lindenbaum;
its world belongs to the first proof set and not the second.
This proves the converse inclusion test. Equality means both
inclusions and hence the theorem biconditional, with its converse
by the same test. Finally membership
\(\beta\in\bigcap|\alpha|_L\) says that each maximal
theory containing \(\alpha\) contains \(\beta\), exactly
the inclusion just proved. If the family is empty this intersection
is the full formula universe; the inclusion test proves every
\(\alpha\to\beta\), as required for inconsistent \(\alpha\).
No finiteness assumption on the initial premises was used.
\end{proof}
\end{annexitem}

\begin{annexitem}{annex:p8.pn.canonical-truth}{thm:p08-pn-canonical-truth}{Existence of the extrema and unrestricted formula induction}
\begin{proof}
If \(|\alpha|_L=|\beta|_L\), the proof-set lemma gives
\(\vdash_L\alpha\leftrightarrow\beta\). RE gives its boxed
biconditional, which belongs to each maximal theory. Thus these
theories agree on the two boxed memberships. If
\(|\alpha|_L\in N_L^{\min}(\Gamma)\), some
\(\Box\beta\in\Gamma\) has the same proof set, so
\(\Box\alpha\in\Gamma\). The converse is its definition.
This proves that the minimum is actually canonical. The maximum
adds no proof sets, so has the same canonical membership test.
Any canonical function must include the minimum and must exclude
exactly the rejected proof sets; every remaining subset is a
non-proof set. It therefore lies below the displayed maximum.
Consistency of L and Lindenbaum applied to the empty premise set
ensure that \(M_L\) is nonempty.

Now induct on \(\alpha\), without requiring its consistency.
Atoms agree by \(V_L(p)=|p|_L\); constants agree by the
unit and empty proof sets. Negation agrees by complement of both
the inductive truth set and the proof set. Conjunction agrees by
intersection of both sets. In the modal case, for every world
\(\Gamma\),
\[
 \mathcal M_L,\Gamma\models\Box\alpha
 \ \Longleftrightarrow\ |\alpha|_L\in N_L(\Gamma)
 \ \Longleftrightarrow\ \Box\alpha\in\Gamma.
\]
This proves its truth set is \(|\Box\alpha|_L\).
Diamond and the other defined connectives follow by their expansions.
For any consistent premise set take a maximal extension; every
member is then true there by the induction. This shows truth in
the \emph{canonical valuation}. It contains no claim about all
arbitrary valuations on the same frame.
\end{proof}
\end{annexitem}

\begin{annexitem}{annex:p8.pn.minimal-not-monotone}{thm:p08-pn-minimal-not-monotone}{An explicit complementary diagonal non-proof set}
\begin{proof}
Set \(\chi=\Box\bottom\) and \(X=|\chi|_{EM}\).
We first construct infinitely many distinct maximal EM theories
outside X, without assuming infinitely many atoms. For each
\(n\geq1\), take a finite relational chain with n edges from
its root to a terminal world, interpreted as principal supplemented
neighborhoods. Give every atom the empty valuation. At the root
\(\Box^k\bottom\) is true precisely when \(k>n\): for
k=1 this is the absence of a successor, and successive boxes
move the same test one link along the chain. The full root theory
\(\Delta_n\) is EM-consistent by EM soundness, decides every
formula, and is therefore maximal. Each omits \(\chi\), while
\(\Delta_n\ne\Delta_m\) for n<m because
\(\Box^{n+1}\bottom\) distinguishes them.

Enumerate all formulas \(\theta_1,\theta_2,\ldots\) and define
\[
 Z=\{\Delta_n:n\geq1,\ \theta_n\notin\Delta_n\}.
\]
It is contained in \(M_{EM}\setminus X\). If
\(Z=|\theta_k|_{EM}\), then \(\Delta_k\in Z\) iff
\(\theta_k\notin\Delta_k\), but also iff
\(\theta_k\in\Delta_k\), a contradiction. Hence Z is
an actual non-proof set. Set \(Y=X\cup Z\). If Y were a
proof set, intersecting it with the proof set
\(|\neg\chi|_{EM}\) would make Z a proof set. Thus Y is
not one and is strictly larger than X.

The one-edge relational chain has \(\Box\Box\bottom\)
true at its root, so its root theory \(\Gamma\) contains
\(\Box\chi\). Consequently
\(X\in N_{EM}^{\min}(\Gamma)\), but Y is not in that
family because every member is a proof set. This is the required
failure of supplementation. It uses only constants and modal
formulas and works with no atoms as well. The non-proof witness
was chosen \emph{in the complement of X}; arbitrary non-proof
sets cannot be substituted for it in the union argument.
\end{proof}
\end{annexitem}

\begin{annexitem}{annex:p8.pn.supplemented-canonical}{thm:p08-pn-supplemented-canonical}{Simultaneous canonical M/C/N conditions}
\begin{proof}
The operation \(\mathcal U^{\rm mon}\) is supplemented:
a witness Y contained in X is also contained in any superset of X.
It contains the original family using each member as its own witness.
Suppose L contains M and
\(|\alpha|_L\in(N_L^{\min}(\Gamma))^{\rm mon}\).
Then some \(\Box\beta\in\Gamma\) has
\(|\beta|_L\subseteq|\alpha|_L\). The proof-set inclusion
test gives theorem \(\beta\to\alpha\). The derived RM
gives theorem \(\Box\beta\to\Box\alpha\); maximal-theory
MP gives \(\Box\alpha\in\Gamma\). Conversely this boxed
membership already puts its proof set in the minimum and hence
the supplement. The supplement is therefore canonical.

If L contains C, members X and Y of the minimum have
\(X=|\alpha|_L,Y=|\beta|_L\) with both boxes in \(\Gamma\).
C and MP put \(\Box(\alpha\land\beta)\) there, so
\(X\cap Y\) is in the minimum. In its supplement,
members A and B have such minimum witnesses X and Y.
Their intersection X\(\cap\)Y is a minimum member and is
contained in A\(\cap\)B, so A\(\cap\)B is in the supplement.
If L contains N, every maximal theory contains
\(\Box\topv\); its proof set is \(M_L\), giving the unit
in the minimum and the supplement. These arguments establish
the combined conditions on one and the same canonical function.
\end{proof}
\end{annexitem}

\begin{annexitem}{annex:p8.pn.six-completeness}{thm:p08-pn-six-completeness}{Soundness and the six full strong-completeness applications}
\begin{proof}
Boolean tautologies are valid on every model. Dual and theorem
RE are sound by extensionality, and MP preserves validity.
The M/C/N correspondences prove soundness of exactly the added
schemas on each displayed class. If a finite premise conjunction
implies \(\varphi\) by a theorem, validity of that implication
gives local semantic consequence. This proves soundness even
when the premise set is infinite, since each derivation uses
only the specified finite subset.

All six logics are consistent: the one-world principal frame
\(N(a)=\{\{a\}\}\) satisfies M, C and N, and no valuation
makes \(\bottom\) true. Soundness thus excludes a bottom theorem
in every selected subsystem. For E take either canonical extreme;
it belongs to the arbitrary-neighborhood class. For EC take the
minimum; the C argument makes all its families intersection-closed.
For EN take the minimum; the N argument inserts its unit. For EM
take the supplemented minimum, which is canonical and supplemented.
For EMC the same choice retains C simultaneously. For EMCN the
same choice also retains the unit, hence has filter families.
Improper filters are allowed: \(\wp(M_L)\) is not excluded.
Each function has the truth lemma and a nonempty world set.

If \(\Gamma\not\vdash_L\varphi\), the set
\(\Gamma\cup\{\neg\varphi\}\) is consistent. Indeed, a
finite inconsistency would prove
\((\delta\land\neg\varphi)\to\bottom\) for a finite
conjunction \(\delta\) of \(\Gamma\), and a classical
tautology would then prove \(\delta\to\varphi\), contrary
to the assumption. A maximal extension is a world of the relevant
canonical model where all of \(\Gamma\) and \(\neg\varphi\)
are true. The chosen model is in precisely the required frame class,
so \(\Gamma\not\models^\ell_{\mathscr F_L}\varphi\).
This proves strong completeness for arbitrary premise sets in
each of the six rows, not just weak completeness. The derived
identity EMCN=K supplies the final source naming bridge.
The identical proof also supplies EMN and ECN using the respective
combined canonical choices when they are used below.
\end{proof}
\end{annexitem}

\begin{annexitem}{annex:p8.pn.augmentation}{thm:p08-pn-augmentation}{Both pointwise conversions and the augmented completeness argument}
\begin{proof}
For \(N^R(w)=\{X:R(w)\subseteq X\}\), every superset of
a member is a member. The set \(R(w)\) itself belongs and is
contained in every member, so \(\bigcap N^R(w)=R(w)\) and
the core belongs. This also works when \(R(w)=\varnothing\):
then the family is the whole powerset. Conversely if N is augmented,
its core \(C_w=\bigcap N(w)\) belongs and supplementation
gives every superset of it. Every member contains it by definition.
Thus \(N(w)=\{X:C_w\subseteq X\}=N^{R_N}(w)\).
In particular the family cannot be empty, since its stipulated
empty-family core W would have to belong.

Induct on formulas with the same valuation on the two pointwise
equivalent frames. Atoms and constants agree; negation and conjunction
preserve agreement. For box the neighborhood clause says exactly
that its truth set contains the relational successor set. Thus the
modal clauses agree, and defined diamond follows. Since the same
argument holds for \emph{every} valuation, frame validity is preserved.

For completeness on augmented frames, construct the relational
canonical K model rather than assuming the supplemented canonical
filter is principal. Put \(W=M_K\), \(V(p)=|p|_K\), and
\(\Gamma R\Delta\) iff
\(B_\Gamma=\{\alpha:\Box\alpha\in\Gamma\}\subseteq\Delta\).
If \(\Box\varphi\in\Gamma\), every successor contains
\(\varphi\). If \(\Box\varphi\notin\Gamma\), the set
\(B_\Gamma\cup\{\neg\varphi\}\) is consistent. Otherwise
finitely many \(\alpha_1,\ldots,\alpha_k\in B_\Gamma\)
would have theorem \((\bigwedge_i\alpha_i)\to\varphi\).
For k>0, the proved positive regular rule and the memberships of
their boxes give \(\Box\varphi\in\Gamma\), contradiction.
For k=0, \(\varphi\) is a theorem and Nec yields the same
contradiction. Lindenbaum supplies a successor containing
\(\neg\varphi\). This proves both modal directions of the
relational truth lemma; atoms and Boolean cases use the maximal-theory
properties exactly as before. A maximal extension of
\(\Gamma_0\cup\{\neg\varphi\}\) therefore witnesses every
failure of local derivability. Turning this relation into \(N^R\)
gives an augmented countermodel. Soundness follows because augmented
families are filters. Together this proves strong completeness on
augmented frames without relying on an unproved relational theorem.

For the distinction of classes, set W=\(\N\) and at every
world take the cofinite filter
\(\mathcal U=\{X\subseteq\N:\N\setminus X\text{ is finite}\}\).
It is supplemented, contains W, and two cofinite sets have cofinite
intersection. For each n the member \(\N\setminus\{n\}\)
omits n; hence its core is empty. That core is not cofinite.
The filter is not augmented and cannot equal the principal family
of any successor set on the same W, since every such family contains
its core. These statements concern pointwise family equality,
not an unproved impossibility of every modal-equivalence construction.
\end{proof}
\end{annexitem}

\begin{annexitem}{annex:p8.pn.regular-relational}{thm:p08-pn-regular-relational}{Full regular canonical truth with possible and impossible worlds}
\begin{proof}
At a normal world the family of all supersets of R(w) is
supplemented and intersection-closed; at an impossible world the
empty family has both closure properties vacuously. Induction
on formulas gives the translation's truth preservation: the modal
clause for normal worlds is membership in that principal family,
and at impossible worlds both presentations make every box false.
The Boolean clauses agree everywhere. The correspondence and RE
proofs establish EMC soundness on these models.

For strong completeness take the stated canonical worlds,
valuation and relation. A world omitting \(\Box\topv\)
contains no \(\Box\alpha\), because RM applied to the theorem
\(\alpha\to\topv\) would imply \(\Box\topv\). Thus at
every impossible world the asserted truth lemma for box holds.
At a normal \(\Gamma\), if \(\Box\varphi\in\Gamma\),
the definition of R and the inner truth induction make \(\varphi\)
true at all successors. For the converse, suppose
\(\Box\varphi\notin\Gamma\). If
\(B_\Gamma\cup\{\neg\varphi\}\) were inconsistent,
some positive finite conjunction of members of \(B_\Gamma\)
would imply \(\varphi\), and the regular rule would force
\(\Box\varphi\in\Gamma\). In the exceptional case where
no member of \(B_\Gamma\) is needed, \(\varphi\) is a
theorem. Apply RM to \(\topv\to\varphi\) and use the
normal-world membership \(\Box\topv\in\Gamma\) to get
the same contradiction. Lindenbaum therefore gives a successor
\(\Delta\supseteq B_\Gamma\cup\{\neg\varphi\}\).
The inner induction makes \(\varphi\) false there. This proves
the modal converse. All atomic and Boolean cases follow from
the valuation and the already proved maximal-theory truth tables.
A maximal extension of any consistent premise set realizes it.
Adjoin the negated conclusion exactly as in the six-system proof
to obtain strong local completeness.

The source's illustrative two-world model is fully specified by
\(W=\{a,b\}\), \(W^N=\{a\}\), \(R=\{(a,b)\}\),
and any atomic valuation. Classical \(\topv\) is true at both
worlds; \(\Box\topv\) is false at b, true at a, and
\(\Box\Box\topv\) false at a. Necessitating an arbitrary
validity therefore fails. Normality is not reinstated merely by
keeping classical Boolean truth at the impossible worlds.
\end{proof}
\end{annexitem}

\begin{annexitem}{annex:p8.pn.filtration}{thm:p08-pn-filtration}{Exact quotient membership through each closure operation}
\begin{proof}
Let \(Q=W/{\sim_\Sigma}\). For every \(\alpha\in\Sigma\),
the original truth set \(T(\alpha)\) is saturated: if w lies
in it and \(w\sim_\Sigma v\), v lies in it. Hence
\(q^{-1}(q[T(\alpha)])=T(\alpha)\). Images of saturated
sets commute with complement, finite intersection and inclusion;
for example, if a quotient class belongs to images of two such
sets, all its representatives belong to both. The same is true
for finite intersections of tracked truth sets even when their
conjunction is not itself in \(\Sigma\).

The generating family \(\mathcal B([w])\) is well defined:
equivalent representatives agree on each \(\Box\alpha\in\Sigma\),
and thus on membership of its truth set in their families. Its
number of distinct generators is at most the number of tracked
box subformulas. If a tracked image \(q[T(\psi)]\) is one
of its generators, it equals \(q[T(\alpha)]\) for an accepted
tracked \(\alpha\). Saturation pulls equality back to
\(T(\psi)=T(\alpha)\), so the original accepts it as well.
This proves exact tracked membership for the finest family.

For C, close the generating family under positive finite intersections.
If a tracked image equals an added intersection
\(\bigcap_i q[T(\alpha_i)]\), pulling back gives
\(T(\psi)=\bigcap_i T(\alpha_i)\). Each factor belonged
to \(N(w)\), so original intersection closure accepts the result.
For N, additionally insert Q. If a tracked image equals Q,
its original saturated truth set is W, which the original unit
condition accepts. Intersecting Q with another generated member
has no additional effect; Q is never substituted for an intersection
of zero generators unless N was selected.

For M, supplement the resulting family. An inserted tracked
\(q[T(\psi)]\) contains either a positive generated
intersection, a single generator when C is absent, or Q when
it was separately inserted by N. Pullback shows that its original
truth set contains the corresponding original accepted set.
Original supplementation therefore accepts \(T(\psi)\).
This proves the converse tracked-membership condition for every
combination of M, C and N. The forward condition holds because
every originally accepted tracked image was a generator from the
start. The constructed families have all the selected conditions:
positive intersections commute with supplementation, as shown in
the canonical closure proof, and the explicitly inserted unit remains.

Now induct over \(\Sigma\). Atomic truth follows from
\(V^f(p)=q[V(p)]\) and saturation for its tracked atoms.
Constants have the stipulated truth sets. Negation and conjunction
follow from the saturated image identities and the inner induction.
For a tracked box the induction gives
\([![\psi]!]_{\mathcal M^f}=q[T(\psi)]\), and the exact
membership condition just proved equates its box truth at [w]
with its original box truth at w. Defined diamond is covered by
first expanding it and keeping all expanded subformulas.

For original proper filters, none of the generated positive
intersections is empty: its pullback belongs to a proper original
filter. Q is nonempty because W is nonempty. Every supplemented
member contains one of these nonempty generators, so the empty
set is never inserted. Finally each quotient world has a unique
truth-value vector on \(\Sigma\), so \(\#Q\leq2^{\#\Sigma}\).
This proves both the truth preservation and the claimed finite bound.
\end{proof}
\end{annexitem}

\begin{annexitem}{annex:p8.pn.finite-canonical}{thm:p08-pn-finite-canonical}{Consistent descriptions, labels and all-valuations non-iterative validity}
\begin{proof}
Since L is consistent, its full maximal world set is nonempty.
Every maximal theory induces a \(\Sigma_\varphi\)-description
by its decisions. That description is consistent, so has its
selected representative. Different descriptions give different
representatives, since their opposite decisions cannot coexist
in one consistent set. Thus \(W_\varphi\) is nonempty and
has at most \(2^{\#\Sigma_\varphi}\) worlds.

If a Boolean combination \(\psi\) of \(\Sigma_\varphi\)
is a theorem, every representative contains it. If it is not a
theorem, \(\{\neg\psi\}\) is consistent and extends to a
maximal theory \(\Delta\). Its description D and the chosen
\(\Gamma_D\) agree on each member of \(\Sigma_\varphi\).
Their Boolean truth tables consequently agree on \(\psi\),
so \(\neg\psi\in\Gamma_D\). This proves the converse
using \emph{L-consistency}, not by treating all descriptions as
independently allowable propositional truth assignments.

The label \(\lambda_D\) belongs to \(\Gamma_D\), and to
no other representative, since each other description has at least
one opposite decision. Therefore \(|\lambda_D|_\varphi\)
is its singleton, and the finite disjunction label has
\(|\lambda_Y|_\varphi=Y\); the empty label is contradictory.
All labels are Boolean combinations of \(\Sigma_\varphi\).
If a Boolean combination \(\psi\) has
\(Y=|\psi|_\varphi\), all representatives contain
\(\lambda_Y\leftrightarrow\psi\); the interface just
proved makes it an L-theorem. RE gives
\(\vdash_L\Box\lambda_Y\leftrightarrow\Box\psi\),
and all representatives contain that equivalence, regardless of
whether the boxed equivalence itself is a Boolean combination
of \(\Sigma_\varphi\).

For every formula on which truth equals membership, negation
and conjunction retain this equality by maximal-theory Boolean
properties. If \(\psi\) is a Boolean combination of
\(\Sigma_\varphi\) and its truth set equals its membership
set Y, then at \(\Gamma_D\), the box clause is
\(Y\in N_\varphi(\Gamma_D)\) iff
\(\Box\lambda_Y\in\Gamma_D\). The last theorem
equivalence makes this iff \(\Box\psi\in\Gamma_D\).
Induction over \(\Sigma_\varphi\) now gives agreement for
every subformula: atomic and constant cases hold by valuation;
Boolean cases use this property; a box has an inner subformula,
which is itself a Boolean combination of the observation set,
so the modal argument applies. Repeated Boolean combination
then gives its entire Boolean closure, and the modal argument
gives boxes of that closure. There is no claim of agreement on
all arbitrary iterated formulas outside these classes.

To prove frame validity of non-iterative axiom schemas, fix an
arbitrary valuation U on \(W_\varphi\). For each atom p
occurring in a non-iterative schema \(\eta\), substitute
\(\lambda_{U(p)}\), and call the resulting formula \(\eta^*\).
Each such label has its correct truth set in the finite canonical
model. Every Boolean modal argument in \(\eta^*\) is a
Boolean combination of these labels and thus of
\(\Sigma_\varphi\); its box has the correct membership set
by the previous paragraph. As \(\eta\) has no modality
inside another, Boolean composition now gives truth equals
membership for \(\eta^*\). It is an L-theorem by schema
substitution, so its membership set is all \(W_\varphi\).
An independent substitution induction equates its truth set in
the canonical valuation with that of \(\eta\) under U:
atoms agree by their labels, Boolean clauses commute, and box
depends only on the inner truth set. Hence \(\eta\) is valid
under U. U was arbitrary. PC, Dual, theorem RE and MP preserve
frame validity, so the frame validates every theorem of L.

If \(\varphi\) is not a theorem, the Boolean interface
gives a representative omitting it, and subformula truth
agreement falsifies it there. This is the required finite
countermodel on a frame validating L. Choosing representatives
uses existence, not a purported procedure for deciding their
consistency. Inconsistent descriptions have never been used.
\end{proof}
\end{annexitem}

\begin{annexitem}{annex:p8.pn.decidability}{thm:p08-pn-decidability}{Finite generator search and the effective non-iterative boundary}
\begin{proof}
Every parameter of the search is finite: j has a finite bound,
each generator is a j-bit subset, each world's list has at most
s entries, and the occurring atoms have finitely many valuations.
Compute the chosen closures by finite set operations on the
finite powerset. For C use positive intersections; for M use
supersets; for N insert the unit. These operations terminate.
The canonical closure argument shows that the generated family
has each specified property simultaneously. Formula evaluation
terminates by recursion down its finite syntax tree.

Every successful candidate is therefore a model in the selected
frame class, so is a genuine satisfiability witness. Conversely,
if the input is satisfiable, filter a satisfying model through
its expanded subformulas. The preceding filtration proof yields
at most \(2^s\) worlds with truth preserved. Its neighborhoods
are obtained by exactly the specified closures from at most s
tracked generators per world. Up to naming the quotient worlds,
this candidate occurs in the finite search, and its valuation
of each occurring atom is enumerated. Thus an unsuccessful
completed search proves unsatisfiability. Searching for the
negation decides validity; the strong-completeness constructions
for each M/C/N combination turn validity into theoremhood.

For a given j, at most \(sj^2\) bits encode all generator
incidences, and at most sj bits encode the occurring atomic
valuations. There are at most \(2^{sj^2+sj}\) such encodings;
summing over \(j\leq2^s\), and allowing finite-power-set
closure and truth evaluation, is bounded by
\(2^{2^{O(s)}}\). For ordinary expanded syntax s is at most
its syntax-tree size, not exponentially large. This is only
a double-exponential exhaustive upper bound; it supplies no
NP or PSPACE lower bound or sharp membership theorem.

For E plus finitely many effective non-iterative schemas, the
finite-canonical theorem gives the same world bound for a
non-theorem's counterframe. Enumerate all neighborhood families
on each domain of that size. For each finite schema pattern,
check every valuation of its finitely many atoms into the
domain's finite powerset; retain a frame exactly when every
pattern is valid under all those valuations. Schema validity
also ensures validity of all formula-substitution instances,
by the same substitution induction, so these retained frames
validate all theorems using E's validity-preserving rules.
The finite counterframe theorem makes this decision procedure
complete. An arbitrary nonrecursive infinite axiom set does
not give a finite effective test of candidates; no decision
claim has been made for that case. Pacuit's cited sharper
complexity results are not replaced by this elementary search.
\end{proof}
\end{annexitem}

\begin{annexitem}{annex:p8.pn.structural-truth}{thm:p08-pn-structural-truth}{Inverse images, every formula case and all disjoint-union conditions}
\begin{proof}
For a bounded morphism f, prove
\(T_{\mathcal M}(\alpha)=f^{-1}(T_{\mathcal M'}(\alpha))\)
by induction. Atomic valuation is its first condition; inverse
images commute with complement relative to the two domains,
with intersection, and with the unit and empty sets, so all
Boolean and constant cases follow. For box the inner induction
and its second condition give
\[
 T_{\mathcal M}(\alpha)\in N(w)
 \ \Longleftrightarrow\ f^{-1}(T_{\mathcal M'}(\alpha))\in N(w)
 \ \Longleftrightarrow\ T_{\mathcal M'}(\alpha)\in N'(f(w)).
\]
That is exactly equivalence of box truth at the two worlds.
Diamond follows by its defined expansion. No step used surjectivity.

The component inclusion \(e_i:W_i\to W\) preserves atomic
valuation, and \(e_i^{-1}(Y)=Y\cap W_i\). The disjoint
union's defining membership clause is precisely the second
bounded-morphism condition for \(e_i\). Thus the induction
preserves truth at all component worlds. If X is accepted at
w in \(W_i\) and X\(\subseteq\)Y, supplementation of
\(N_i(w)\) accepts \(Y\cap W_i\), so the union accepts Y.
If X,Y are accepted, intersection closure of the component
accepts \((X\cap W_i)\cap(Y\cap W_i)
=(X\cap Y)\cap W_i\), so the union accepts their intersection.
Its unit intersects \(W_i\) in the component unit, which is
accepted when each component has N. These are separate proofs
of the three preserved conditions.
\end{proof}
\end{annexitem}

\begin{annexitem}{annex:p8.pn.one-step}{thm:p08-pn-one-step}{Correct root countermodels, theorem-level matrices and intersection folding}
\begin{proof}
If the atomic lists overlap, the sequent is valid propositionally.
For EM, a globally valid matrix implication gives truth-set
inclusion in every model of the class; supplementation turns an
accepted antecedent truth set into an accepted consequent one.
For E use equality of truth sets from the matrix biconditional.
For EC, if a nonempty conjunction of matrix antecedents is
globally equivalent to a matrix consequent, intersection closure
accepts that conjunction's truth set whenever all the corresponding
boxes hold, and equality accepts the consequent. These prove
all sufficient conditions.

We construct the converses without assigning root neighborhoods
in terms of truth values that depend on those same neighborhoods.
Assume the atomic lists are disjoint. For EM, if none of the
matrix implications is valid, choose for each pair i,j a pointed
monotonic model satisfying \(\varphi_i\land\neg\psi_j\).
Rename worlds to make these models disjoint and take their source
disjoint union \(\mathcal M_0\) on \(W_0\).
If there are no such pairs, adjoin a one-world empty-family
component instead, to keep \(W_0\) nonempty. Let
\(A_i=T_{\mathcal M_0}(\varphi_i)\) and
\(B_j=T_{\mathcal M_0}(\psi_j)\). The witnesses give
\(A_i\not\subseteq B_j\) for all pairs.

Add a fresh root r. For old w in \(W_0\), define acceptance
of \(X\subseteq W_0\cup\{r\}\) iff
\(X\cap W_0\in N_0(w)\). Thus the old component embedding
is a bounded morphism; no old formula truth changes. At r make
exactly the antecedent atomic list true (other atoms are false).
Put \(N(r)=\{X:\exists i,\ A_i\subseteq X\}\).
The whole model is monotonic. Its full antecedent truth set
restricts to \(A_i\) on old worlds, and contains \(A_i\),
so every \(\Box\varphi_i\) is true at r. Its full consequent
truth set restricts to \(B_j\); it cannot contain any
\(A_i\), because the corresponding old witness persists.
Hence every \(\Box\psi_j\) is false at r. The disjoint
atomic lists are true and false respectively at r. This falsifies
the sequent, including the cases with one of the box lists empty.

For E, choose for each pair a pointed arbitrary neighborhood model
where \(\varphi_i\leftrightarrow\psi_j\) is false. Its
direction of inequality may differ from pair to pair. Their disjoint
union gives \(A_i\ne B_j\). Make the same old-world extension
and atomic root valuation, but set
\[
 N(r)=\{A_i,A_i\cup\{r\}:1\leq i\leq k\}.
\]
Regardless of its truth at r, the full truth set of \(\varphi_i\)
is one of these two sets. The full truth set of \(\psi_j\)
cannot equal either, because its restriction \(B_j\) differs
from every \(A_i\). Thus precisely the required antecedent
boxes hold and consequent boxes fail. This proves the E converse.

For EC, assume that no nonempty conjunction of any selected
antecedent matrices is globally equivalent to any \(\psi_j\)
on the intersection-closed class. For each nonempty
\(I\subseteq\{1,\ldots,k\}\) and j choose a pointed
intersection-closed countermodel of that biconditional. Their
disjoint union gives, for every such I,j,
\(A_I=\bigcap_{i\in I}A_i\ne B_j\). Use the same old-world
extension, and at r take the positive finite intersection closure
of \(\{A_i,A_i\cup\{r\}:1\leq i\leq k\}\).
Every resulting member has restriction \(A_I\) for a nonempty
I. Therefore none equals the full truth set of a consequent;
every antecedent truth set is one of the original generators.
The root and old-world families are intersection-closed. The same
atomic valuation falsifies the sequent. If k=0, take the empty
root family; no zero-fold intersection is inserted. This proves
the exact EC characterization, including its singleton case.

For the valid EC folding inference, any world satisfying
\(\Box\alpha\) and \(\Box\beta\) satisfies
\(\Box(\alpha\land\beta)\). Thus if the sequent with the
boxed conjunction and any unchanged contexts is already valid,
the sequent with the two separate boxed antecedents is valid.
The opposite semantic projection is M and is unavailable in EC;
the C+N/not-M frame in the finite separations witnesses that failure.

The printed contextual EM and E inferences cannot be justified
with arbitrary unchanged \(\Gamma,\Delta\). On the earlier
unit two-world countermodel with p true at a only, at a we have
\(p\), \(\Box\topv\) and \(\neg\Box p\).
The premise sequent \(p,\topv\Rightarrow p\) is globally
valid, but the proposed conclusion
\(p,\Box\topv\Rightarrow\Box p\) is false at a.
For the two-premise E version, both
\(p,\topv\Rightarrow p\) and \(p,p\Rightarrow\topv\)
are valid, yet the same conclusion fails. This model is also
monotonic. Hence both counterexamples fall within the advertised
classes. Sound replacement or monotonicity requires a separately
valid matrix equivalence or implication, after which ordinary
context weakening is harmless; a context-dependent implication
is not such a theorem. The repaired root construction above
does not use the source's unjustified assertion that the root's
principal family alone forces all listed boxes.
\end{proof}
\end{annexitem}

\begin{annexitem}{annex:p8.pn.general-completeness}{thm:p08-pn-general-completeness}{Admissible substitutions and the complete local argument}
\begin{proof}
The proof-set identities make \(P_L\) a Boolean algebra with
unit \(M_L\), empty set, complements and finite intersections.
Canonical membership gives
\(m_{N_L^{\min}}(|\alpha|_L)=|\Box\alpha|_L\), so it
is closed under its required modal operation. Thus the canonical
triple is a well-defined general frame. In any general model,
induction keeps all formula truth sets in the admissible algebra:
atoms lie there by admissibility, constants by its unit and empty
set, Boolean clauses by closure, and box by \(m_N\).

Fix an arbitrary admissible valuation U on the canonical general
frame, and a theorem \(\eta\) of L. Each of its finitely
many occurring atoms p has \(U(p)=|\rho_p|_L\) for some
formula \(\rho_p\). Choose these formulas and substitute them
simultaneously into \(\eta\), leaving other atoms fixed.
For any subformula \(\alpha\) of \(\eta\), induction proves
\[
 [![\alpha]!]_{(M_L,N_L^{\min},U)}
   =[![\alpha[\rho_p/p]]!]_{(M_L,N_L^{\min},V_L)}.
\]
For an atom this is the canonical truth lemma applied to its
chosen proof set. For constants the truth sets agree directly.
For negation and conjunction the equality is preserved by
complement and intersection. For box, the induction equates its
two argument truth sets, and both models use the same function,
so the sets of worlds accepting them are equal. This proves the
substitution identity at \(\eta\).

Substitution closure of L makes \(\eta[\rho_p/p]\) an
L-theorem. It belongs to every maximal theory and hence has
canonical truth set \(M_L\). The identity makes \(\eta\)
true everywhere under U. Since U was arbitrary, the general
frame validates L. Consistency was necessary to keep its
canonical world set nonempty. An arbitrary extension not closed
under formula substitution would not justify the theorem step.

The canonical valuation is admissible. Given
\(\Gamma\not\vdash_L\varphi\), take the maximal extension
of \(\Gamma\cup\{\neg\varphi\}\) constructed earlier.
The canonical truth lemma gives a local countermodel on this
L-valid general frame. This proves strong completeness for the
class of all L-general frames. Soundness holds because every
theorem is valid by the class's definition, and a theorem finite
premise implication preserves truth at a world. No arbitrary
valuation outside \(P_L\) has been inserted into this argument,
so it proves no unrestricted ordinary-frame completeness claim.
\end{proof}
\end{annexitem}
